\section{Conclusioni}

Il progetto A.L.I.C.E. realizza una piattaforma web semplice per il monitoraggio remoto di pazienti
anziani. La soluzione mette insieme tre aree concrete: farmaci, esercizi fisici e umore quotidiano.
L'operatore usa una dashboard per controllare i pazienti, mentre il paziente usa pagine lineari con
azioni brevi e leggibili.

\subsection{Risultati raggiunti}

\begin{itemize}
    \item \textbf{Separazione dei ruoli:} operatore e paziente hanno pagine, login e percorsi distinti.
    \item \textbf{Dashboard operatore:} l'operatore visualizza pazienti, stato, metriche e dettaglio individuale.
    \item \textbf{Area paziente semplificata:} il paziente accede a umore, esercizi e farmaci partendo da una home con tre scelte principali.
    \item \textbf{Database coerente:} lo schema contiene tabelle dedicate a profili, operatori, pazienti, farmaci, esercizi e log.
    \item \textbf{Navigazione protetta:} il proxy Next.js gestisce i redirect principali in base alla sessione Supabase.
    \item \textbf{Suggerimento AI:} l'operatore puo' ricevere una frase sintetica di supporto basata sui dati del paziente.
\end{itemize}

\subsection{Benefici attesi}

\paragraph{Per il paziente}
\begin{itemize}
    \item interfaccia piu' semplice rispetto a un gestionale sanitario tradizionale;
    \item meno passaggi per registrare umore, farmaci ed esercizi;
    \item conferma immediata delle azioni completate.
\end{itemize}

\paragraph{Per l'operatore}
\begin{itemize}
    \item visione centralizzata dei pazienti seguiti;
    \item controllo piu' rapido di aderenza farmacologica, esercizi e umore;
    \item possibilita' di modificare piano e dati del paziente da una sola pagina.
\end{itemize}

\subsection{Limiti attuali}

Il progetto copre il flusso principale, ma resta una piattaforma didattica. Alcune funzioni sono
volutamente semplici: non sono presenti notifiche automatiche, videochiamate, app mobile nativa o
integrazione con dispositivi medici. Questa scelta mantiene il codice piu' leggibile e adatto alla
presentazione universitaria.

\subsection{Sviluppi futuri}

Possibili miglioramenti futuri sono:

\begin{itemize}
    \item notifiche per ricordare al paziente farmaci ed esercizi;
    \item grafici piu' avanzati per lo storico dell'operatore;
    \item esportazione dei report del paziente;
    \item videochiamata semplificata tra operatore e paziente;
    \item app mobile o Progressive Web App installabile.
\end{itemize}

\newpage
